%!TEX root=paper.tex
\section{Demonstration scenario}

We will demonstrate CLUE with an open-source dbt project, Jaffle Shop
~\cite{jaffleShop}, which contains six source tables: customers, orders, order
items, products, supplies, and stores. The demonstration guides users through 11
steps (as annotated in Figure 2). We impersonate Sheldon, introduced in Example
1, who has recently joined Jaffle Shop as a data analyst. While exploring the
company's operations dashboard, he notices that the PLTV metric is lower than
expected. He wants to understand why this is happening and explore how the
company could increase the metric by 20\% by the end of the year.

We will demonstrate CLUE with an open-source dbt project, Jaffle Shop
~\footnote{Jaffle Shop dataset: \url{https://github.com/dbt-labs/jaffle_shop}},
which contains six source tables: customers, orders, order items, products,
supplies, and stores. The demonstration guides users through 11 steps (as
annotated in Figure 2). We impersonate Sheldon, introduced in Example 1, who has
recently joined Jaffle Shop as a data analyst. While exploring the company's
operations dashboard, he notices that the PLTV metric is lower than expected. He
wants to understand why this is happening and explore how the company could
increase the metric by 20\% by the end of the year.

\textbf{Step \stepCounter{A}, \stepCounter{B}, \& \stepCounter{C} (Searching for
a metric and viewing the description)} Sheldon first searches for the metric
name in CLUE. The left panel displays a snapshot of the metric along with an
AI-generated description explaining how the metric is derived.

\textbf{Step \stepCounter{D} (Viewing metric provenance)} On the top right
panel, CLUE presents the metric-centric lineage of the searched metric as a tree
structure. Each child node represents a submetric of the parent node and is
marked with either \texttt{\#} or \texttt{\%} so users can easily distinguish
their types. Mathematical operations between sibling nodes are annotated with a
yellow diamond on the connecting branches. For example, the root node and the
nodes at level 1 represent the equation $\metric{PLTV} = \metric{Probability of
Active} \times \metric{Expected Number of Orders} \times \metric{Expected Order
Value}$. To avoid overwhelming users, only the first-level submetrics are
displayed by default. Users can click a node to expand the next level.

\textbf{Step \stepCounter{E} (Exploring data provenance)} To ensure the
calculation is correct, Sheldon invites the metric owner to review the data
provenance of the entire metric path. By right-clicking a leaf node and clicking
``\textit{Open in Transformation View}'', CLUE extracts and visualizes the data
flow and transformation steps as a DAG in the panel below. Unlike existing
lineage tools, CLUE allows users to view in-line code context directly within
the lineage graph by hovering over the symbol representing each operation. CLUE
also supports exploring transformations across multiple paths by allowing users
to switch between tabs.

\textbf{Step \stepCounter{F}, \stepCounter{G}, \& \stepCounter{H} (Simulating
metric causality through top-down reasoning)} After confirming that the metric
is correctly defined, Sheldon wants to explore how to plan for the 20\% growth
target in PLTV. She clicks ``\textit{Launch Metric Simulator}'' in
\stepCounter{F} and sets the target growth to 20\% in \stepCounter{G}. CLUE then
initializes sliders for each first-level metric. When the user adjusts any
slider, the remaining sliders automatically adjust to maintain the target value.

\textbf{Step \stepCounter{H}, \stepCounter{I} \&  \stepCounter{J} (Requesting
suggestions)} The user clicks ``\textit{Recommend strategies}'' in
\stepCounter{H}, and CLUE suggests several approaches to achieve the goal set in
\stepCounter{G}. The data change summary of an option will show in
\stepCounter{J} when the user selects it. 

\textbf{Step \stepCounter{K} \& \stepCounter{J} (Simulating value changes of the
selected recommendation and viewing the result)} 

To answer the question ``\textit{How much would PLTV change if we obtained more
orders from new visitors?}'', Sheldon selects the second suggestion and clicks
``\textit{Start simulation'}' in \stepCounter{K}. CLUE then computes the updated
values for each metric and displays the results in \stepCounter{K}. Changes are
highlighted in green or red to indicate increases or decreases in metric values.

\emph{\underline{Demonstration engagement.}} Our target users include
non-technical decision makers and data professionals who need to understand how
a complex metric is derived. After the guided demonstration, users can explore
the causal relationships between a final metric and its component metrics
through clear visualizations, and observe how changes in underlying data
propagate to the metric value.